\documentclass[11pt]{article}
\usepackage[a4paper,margin=1in]{geometry}
\usepackage{amsmath}
\usepackage{booktabs}

\title{Low-Rank Adapters for Small Sentiment Classifiers: A Synthetic Demo}
\author{Aster Research Demo}
\date{October 2026}

\begin{document}
\maketitle

\begin{abstract}
This manuscript belongs to the Aster harness demo. Every run in it was produced by a
seeded simulator, not by training a model, so the numbers exercise the workflow
rather than support a scientific claim.
\end{abstract}

\section{Question}
How much validation accuracy does a LoRA adapter \cite{hu2022lora} gain from a higher
rank when a small encoder \cite{sanh2019distilbert} is fine-tuned on a five-class
sentiment task \cite{socher2013sst}, and which learning rate keeps training stable?

\section{Method}
All runs use AdamW \cite{loshchilov2019adamw}, batch size 32 and 600 optimizer steps.
Each run was launched with the research harness, which records the code at launch,
the exact command and every metric event.

\section{Results}
Table~\ref{tab:rank} reports validation accuracy after 600 steps.
With rank~16 the adapter reaches 82.9\% validation accuracy.
Rank~64 raises it to 86.4\% on the same split.
Moving from rank~16 to rank~64 therefore adds about 2.9 points.

\begin{table}[h]
  \centering
  \begin{tabular}{lcc}
    \toprule
    Rank & Trainable parameters & Val.\ accuracy \\
    \midrule
    4  & 0.3M & 80.0\% \\
    16 & 1.2M & 82.9\% \\
    64 & 4.7M & 86.4\% \\
    \bottomrule
  \end{tabular}
  \caption{LoRA rank ablation (learning rate $3\times10^{-4}$, seed 2).}
  \label{tab:rank}
\end{table}

In the learning-rate sweep, $3\times10^{-4}$ reaches 83.2\% validation accuracy, while
$10^{-3}$ diverges at step 300.
At $10^{-4}$ the validation loss is still 1.047 after 600 steps.

\section{Limitations}
A single seed per configuration cannot separate the rank effect from run-to-run noise.

\bibliographystyle{plain}
\bibliography{library}

\end{document}
